%!TEX root = TheoreticalPhysicsC.tex

\section{Lecture: The Finite Potential Well}
We now consider a particle in a potential well of finite depth $V_0$ and width $L$:
\begin{equation}
	V(x) = \begin{cases}
		-V_0, \quad &|x|\le \frac{L}{2}\\
		0, \quad &|x| > \frac{L}{2}.
	\end{cases}
\end{equation}
We shall take up the following questions: (a) Are there bound states (i.e.\ energy eigenvectors with eigenenergy $E\le 0$)? (b) What differences from the infinitely high potential well do we observe?

\begin{intuition}
	If we fire a wave packet at the well from outside, then part of it flows through, part is reflected, and a remainder stays sitting in the well. A scattering process, then, in principle the same one as at the LHC, where the bound states are new particles.
\end{intuition}

The full dynamics $U_t$ mixes bound and free states and is too complicated; a bound state, however, does not move in time and is therefore an eigenvector of $\widehat{H}$.

In order to obtain the solution of the eigenvector equation  
\begin{equation}
	\widehat{H}|\Psi\rangle = E|\Psi\rangle
\end{equation}
it is helpful to define the constants
\begin{equation}
	\kappa = \frac{\sqrt{2m|E|}}{\hbar}\quad \text{and}\quad q = \frac{\sqrt{2m(E+V_0)}}{\hbar}
\end{equation}


We obtain the associated wave functions from
\begin{equation}
	\langle x|\widehat{H}|\Psi\rangle = \left(-\frac{\hbar^2}{2m}\frac{d^2}{dx^2} + V(x)\right)\Psi(x) = E\Psi(x) = E\langle x|\Psi\rangle.
\end{equation}
After rearranging and substituting we obtain
\begin{equation}
	\begin{cases}
		\frac{d^2}{dx^2}\Psi(x) = -q^2 \Psi(x), \quad &\text{for $|x|<\frac{L}{2}$}\\
		\frac{d^2}{dx^2}\Psi(x) = \kappa^2 \Psi(x), \quad &\text{for $|x|>\frac{L}{2}$}.
	\end{cases}
\end{equation}
By symmetry the eigenfunctions may be chosen to be even or odd. Even means: replace $x$ by $-x$ and the same function comes back. We restrict ourselves to $E<0$; the infinitely many solutions with $E>0$ are \emph{scattering states}, plane waves outside, and are left to the exercises. The even solutions read
\begin{equation}
	\Psi(x) = \begin{cases}
		A_1 e^{\kappa x}, \quad & x< -\frac{L}{2} \\
		A_2 \cos(qx), \quad & |x| \le \frac{L}{2} \\
		A_1 e^{-\kappa x}, \quad & x> \frac{L}{2}
	\end{cases}
\end{equation}
and the odd solutions
\begin{equation}
	\Psi(x) = \begin{cases}
		-B_1 e^{\kappa x}, \quad & x< -\frac{L}{2} \\
		B_2 \sin(qx), \quad & |x| \le \frac{L}{2} \\
		B_1 e^{-\kappa x}, \quad & x> \frac{L}{2}.
	\end{cases}
\end{equation}
Outside they fall off exponentially, but they do not vanish: the particle is to be found with small probability in the classically forbidden region, and it does so in a state that does not change in time.

All wave functions are assumed to be \emph{continuous} and \emph{differentiable}. This yields matching conditions at $x = \pm L/2$. The justification is physical, not mathematical: a jump in $\Psi$ means an infinite first derivative, and that is tied to the momentum operator -- the particle would have infinite momentum there, and infinite kinetic energy. Impossible this is not -- such systems exist -- but we demand finite kinetic energy. That is precisely where the quantisation lies. For the even solutions we have
\begin{enumerate}
	\item continuity: $A_1 e^{-\frac{\kappa L}{2}} = A_2 \cos\left(\frac{qL}{2}\right)$,
	\item differentiability: $-A_1\kappa e^{-\frac{\kappa L}{2}} = -A_2 q\sin\left(\frac{qL}{2}\right)$.
\end{enumerate}
Dividing the second condition by the first, we obtain
\begin{equation}\label{eq:symmcond}
	\kappa = q\tan\left(\frac{qL}{2}\right).
\end{equation}
For the odd solutions the conditions read  
\begin{enumerate}
	\item continuity: $B_1 e^{-\frac{\kappa L}{2}} = B_2 \sin\left(\frac{qL}{2}\right)$,
	\item differentiability: $-B_1\kappa e^{-\frac{\kappa L}{2}} = B_2 q\cos\left(\frac{qL}{2}\right)$.
\end{enumerate}
In the same way one is led to the conclusion:
\begin{equation}\label{eq:unsymm}
	\kappa = -q\cot\left(\frac{qL}{2}\right).
\end{equation}

Let us now consider condition (\ref{eq:symmcond}) for the even solutions: squaring both sides gives
\begin{equation}
	\kappa^2 = q^2\tan^2\left(\frac{qL}{2}\right).
\end{equation}
By the definition of $\kappa$ we obtain
\begin{equation}
	\kappa^2 = \frac{2m|E|}{\hbar^2} = \frac{2mV_0}{\hbar^2} - \frac{2m(V_0+E)}{\hbar^2} = \frac{2mV_0}{\hbar^2} - q^2.
\end{equation}
If we introduce the constants
\begin{equation}
	u\equiv \frac{qL}{2}\quad \text{and}\quad \alpha^2 \equiv \frac{m V_0 L^2}{2\hbar^2}
\end{equation}
we obtain
\begin{equation}
	\alpha^2-u^2 = u^2\tan^2(u).
\end{equation}
After rearranging:
\begin{equation}
	\frac{\alpha^2}{u^2} = 1+\tan^2(u) = \frac{1}{\cos^2(u)}.
\end{equation}
so that
\begin{equation}
	|\cos(u)| = \frac{u}{\alpha}.
\end{equation}
This \emph{transcendental} equation is solved graphically. Because of the conditions $\kappa>0$, $q>0$ and $\kappa = q\tan(u)$ we must have $\tan(u)>0$, so that only solutions in $0<u<\frac{\pi}{2}$ and $\pi < u < \frac{3\pi}{2}$ are admissible. To sum up: there is at least one bound state, and altogether finitely many states with an even wave function.

\begin{merksatz}
	Everything turns on $\alpha$ alone: small for a light particle, a shallow or a narrow well, large for a heavy particle, a deep or a wide well. The larger $\alpha$, the more bound states -- an even one is always there.
\end{merksatz}


Condition (\ref{eq:unsymm}) for the odd solutions is treated in the same way: in this case squaring both sides gives
\begin{equation}
	\kappa^2 = q^2\cot^2\left(u\right).
\end{equation}
Hence we obtain
\begin{equation}
	\alpha^2-u^2 = u^2\cot^2(u).
\end{equation}
After rearranging:
\begin{equation}
	\frac{\alpha^2}{u^2} = 1+\cot^2(u) = \frac{1}{\sin^2(u)}.
\end{equation}
so that
\begin{equation}
	|\sin(u)| = \frac{u}{\alpha}.
\end{equation}
This equation is likewise solved graphically: because $\kappa>0$, $q>0$ and $\kappa = -q\cot(u)$ we must have $\cot(u)<0$, so that only solutions in $\frac{\pi}{2}<u< \pi$ and $\frac{3\pi}{2} < u < 2\pi$ are possible. Unlike $|\cos(u)|$, $|\sin(u)|$ starts at zero at the origin; for small $\alpha$ there can therefore be no odd solution at all. The intersection at $u=0$ does not count: there we would have $q=0$ and the sine solution would vanish identically -- the usual trap when solving graphically.

% fig_vl13_topf.tex — Endlich hoher Potentialtopf: die gebundenen
% Zustaende mit ihren exponentiellen Auslaeufern, und die graphische
% Loesung, die abzaehlt, wie viele es davon gibt.
%
% GEOMETRY FROM THE FORMULA.  Nothing here is drawn by eye.  With the
% abbreviations of the lecture,
%
%       u = qL/2 ,   alpha^2 = m V_0 L^2 / (2 hbar^2) ,
%
% the matching conditions kappa = q tan(u) and kappa = -q cot(u) reduce
% to the two transcendental equations
%
%       |cos u| = u/alpha   (even)      |sin u| = u/alpha   (odd),
%
% each restricted to the branch on which the sign condition holds:
% tan u > 0 on (0,pi/2) and (pi,3pi/2), cot u < 0 on (pi/2,pi) and
% (3pi/2,2pi).  Those four intervals are exactly the four DESCENDING
% half-arches of |cos u| and |sin u|, alternating even/odd — which is
% why panel (b) can be read off as "count the arches the straight line
% u/alpha still reaches".  The ascending halves are the branches the
% sign conditions throw away, the crossing at u = 0 among them.
%
% For alpha = 4 there are ceil(2 alpha/pi) = 3 bound states, with roots
%
%       u_1 = 1.25235   u_2 = 2.47458   u_3 = 3.59530 ,
%
% and with v = kappa L/2 = sqrt(alpha^2 - u^2),
%
%       v_1 = 3.79890   v_2 = 3.14269   v_3 = 1.75322 .
%
% Panel (a) then follows without any further choice:
%
%       E_n / V_0 = u_n^2/alpha^2 - 1  = -0.90198, -0.61728, -0.19211 ,
%       psi_n(x)  = cos(2 u_n x/L)  or  sin(2 u_n x/L)      inside,
%                   (+-) (u_n/alpha) exp(-2 v_n (|x|/L - 1/2))  outside,
%
% the amplitude at the wall being u_n/alpha = |cos u_n| resp. |sin u_n|
% — the transcendental equation itself.  So the tail of the topmost
% state starts almost three times as high AND decays 2.2x more slowly
% than the ground state's (1/kappa_3 = 0.285 L against 1/kappa_1 =
% 0.132 L): the weaker the binding, the further the particle leaks out.
% That comparison is what the picture is for; prose cannot do it.
%
% DESIGN DECISIONS
%   * colour encodes the quantum number n and nothing else, and it is
%     the SAME code in both panels, so the arch a state comes from in
%     (b) can be found again as a wavefunction in (a).  Entries 1, 2, 4,
%     5 of the ordered `tpclevels` ramp (navy, teal, orange, red);
%     entry 3, munsell, is skipped on purpose — beside cgblue it is the
%     same teal to the eye, and arches 2 and 3 in (b) are neighbours;
%   * the fourth colour appears only in (b): arch 4 is the state that
%     alpha = 4 does not yet support.  It is drawn, not omitted, because
%     the empty arch is the point — increase alpha, the line flattens,
%     and the fourth state appears;
%   * everything that is reference rather than data is neutral grey: the
%     energy levels, the -V_0 guide, the straight line u/alpha, and the
%     dotted ascending half-arches;
%   * classically forbidden region hatched, not filled, so the leakage
%     of psi into it survives greyscale printing;
%   * amplitude 0.12 V_0 in (a) is the largest that keeps psi_2 clear of
%     E_1, the tightest gap (E_2 - E_1 = 0.285 V_0).

\begin{center}
\begin{tikzpicture}

% Roots of |cos u| = u/alpha and |sin u| = u/alpha at alpha = 4, and the
% quantities that follow from them.  Local to this tikzpicture.
\def\uA{1.25235}\def\vA{3.79890}\def\EA{-0.90198}\def\bA{0.31309}
\def\uB{2.47458}\def\vB{3.14269}\def\EB{-0.61728}\def\bB{0.61864}
\def\uC{3.59530}\def\vC{1.75322}\def\EC{-0.19211}\def\bC{-0.89883}
\def\amp{0.12}

% ─────────────────────────────────────────────────────────────
%  (a)  The three bound states, each on its own level
% ─────────────────────────────────────────────────────────────
\begin{axis}[
    name=well,
    tpcplot,
    width=7.6cm, height=6.6cm,
    axis lines=none,
    xmin=-1.30, xmax=1.30,
    ymin=-1.24, ymax=0.24,
    clip=false,
    title={(a) \tpcifenglish{Bound states, $\alpha=4$}%
                            {Gebundene Zust\"ande, $\alpha=4$}},
    domain=-1.24:1.24, samples=401,
]

  % Classically forbidden for every bound state: outside the well V = 0
  % lies above E < 0.  Hatched, so psi can be drawn straight over it.
  \addplot[forbidden, draw=none, domain=-1.24:-0.5, samples=2] {-1.24} \closedcycle;
  \addplot[forbidden, draw=none, domain=0.5:1.24,   samples=2] {-1.24} \closedcycle;

  % The well itself
  \draw[potential] (axis cs:-1.24,0)  -- (axis cs:-0.5,0);
  \draw[potential] (axis cs:-0.5,0)   -- (axis cs:-0.5,-1);
  \draw[potential] (axis cs:-0.5,-1)  -- (axis cs:0.5,-1);
  \draw[potential] (axis cs:0.5,-1)   -- (axis cs:0.5,0);
  \draw[potential, -{Stealth[length=5pt]}] (axis cs:0.5,0) -- (axis cs:1.30,0);

  % Energy levels E_n/V_0 = u_n^2/alpha^2 - 1, and the floor -V_0.
  % Reference rules, hence grey: colour is already carrying n.
  \addplot[levellinegrey, domain=-1.24:1.24, samples=2] {\EA};
  \addplot[levellinegrey, domain=-1.24:1.24, samples=2] {\EB};
  \addplot[levellinegrey, domain=-1.24:1.24, samples=2] {\EC};
  \addplot[levellinegrey, domain=-1.24:-0.5, samples=2] {-1};

  % Eigenfunctions about their levels.  Inside: cos/sin(2 u_n x/L).
  % Outside: amplitude at the wall +-u_n/alpha, times exp(-2 v_n (|x|/L - 1/2)).
  \addplot[wavefn, spacecadet] {\EA + \amp*(
       (abs(x)<0.5)*cos(deg(2*\uA*x))
     + (abs(x)>=0.5)*\bA*exp(-2*\vA*(abs(x)-0.5)) )};
  \addplot[wavefn, cgblue]     {\EB + \amp*(
       (abs(x)<0.5)*sin(deg(2*\uB*x))
     + (abs(x)>=0.5)*((x>0)*2-1)*\bB*exp(-2*\vB*(abs(x)-0.5)) )};
  \addplot[wavefn, birrenorange] {\EC + \amp*(
       (abs(x)<0.5)*cos(deg(2*\uC*x))
     + (abs(x)>=0.5)*\bC*exp(-2*\vC*(abs(x)-0.5)) )};

  % Level labels
  \node[lbl, anchor=east] at (axis cs:-1.30,\EA) {$E_{1}$};
  \node[lbl, anchor=east] at (axis cs:-1.30,\EB) {$E_{2}$};
  \node[lbl, anchor=east] at (axis cs:-1.30,\EC) {$E_{3}$};
  \node[lbl, anchor=east] at (axis cs:-1.30,-1)  {$-V_{0}$};

  % Axis furniture
  \node[lbl, anchor=west]       at (axis cs:1.32,0)     {$x$};
  \node[lbl, anchor=south east] at (axis cs:1.24,0.03)  {$V=0$};
  \node[lbl, anchor=south]      at (axis cs:-0.5,0.04)  {$-\tfrac{L}{2}$};
  \node[lbl, anchor=south]      at (axis cs:0.5,0.04)   {$\tfrac{L}{2}$};

  % Well width
  \draw[{Stealth[length=4pt]}-{Stealth[length=4pt]}, birrengray, thin]
    (axis cs:-0.5,-1.14) -- (axis cs:0.5,-1.14)
    node[midway, lbl, fill=white, inner sep=1pt] {$L$};

  % Penetration depth 1/kappa_n = L/(2 v_n): 0.132 L for the ground state
  % against 0.285 L for the topmost one.  White underlay so the measure
  % reads over the hatching.
  \draw[white, line width=2.6pt] (axis cs:0.5,-0.96) -- (axis cs:0.632,-0.96);
  \draw[white, line width=2.6pt] (axis cs:0.5,-0.40) -- (axis cs:0.785,-0.40);
  \draw[{Stealth[length=3pt]}-{Stealth[length=3pt]}, birrengray!60!black, thin]
    (axis cs:0.5,-0.96) -- (axis cs:0.632,-0.96);
  \draw[{Stealth[length=3pt]}-{Stealth[length=3pt]}, birrengray!60!black, thin]
    (axis cs:0.5,-0.40) -- (axis cs:0.785,-0.40);
  \node[lbl, font=\scriptsize, anchor=west, fill=white, inner sep=1pt]
    at (axis cs:0.66,-0.96) {$1/\kappa_{1}$};
  \node[lbl, font=\scriptsize, anchor=west, fill=white, inner sep=1pt]
    at (axis cs:0.82,-0.40) {$1/\kappa_{3}$};

  % The hatched zone, named once
  \node[lbl, text=birrengray, font=\scriptsize, align=center,
        fill=white, inner sep=1pt]
    at (axis cs:-0.87,-1.13)
    {\tpcifenglish{classically\\forbidden}{klassisch\\verboten}};
\end{axis}

% ─────────────────────────────────────────────────────────────
%  (b)  |cos u| = u/alpha and |sin u| = u/alpha: counting the states
% ─────────────────────────────────────────────────────────────
\begin{axis}[
    tpcplot,
    at={(well.south east)}, anchor=south west, xshift=12mm,
    width=5.9cm, height=6.6cm,
    xlabel={$u=\tfrac{qL}{2}$},
    xmin=0, xmax=5.45, ymin=0, ymax=1.55,
    xtick={0,1.5707963,3.1415927,4.712389},
    xticklabels={$0$,$\tfrac{\pi}{2}$,$\pi$,$\tfrac{3\pi}{2}$},
    ytick={0,0.5,1},
    grid=major,
    samples=120,
    title={(b) \tpcifenglish{Graphical solution}{Grafische L\"osung}},
]

  % Branches killed by the sign conditions tan u > 0 / cot u < 0: the
  % ascending half-arches.  Dotted grey — they are not solutions.
  \addplot[birrengray!75, line width=0.7pt, densely dotted,
           domain=0:1.5707963]          {abs(sin(deg(x)))};
  \addplot[birrengray!75, line width=0.7pt, densely dotted,
           domain=1.5707963:3.1415927]  {abs(cos(deg(x)))};
  \addplot[birrengray!75, line width=0.7pt, densely dotted,
           domain=3.1415927:4.712389]   {abs(sin(deg(x)))};
  \addplot[birrengray!75, line width=0.7pt, densely dotted,
           domain=4.712389:5.45]        {abs(cos(deg(x)))};

  % The straight line u/alpha.  Reference, not data: grey.
  \addplot[birrengray!55!black, line width=0.9pt,
           domain=0:5.45, samples=2] {x/4};

  % The four admissible branches, coloured by the state they carry.
  \addplot[spacecadet,   line width=1.0pt, domain=0:1.5707963]         {abs(cos(deg(x)))};
  \addplot[cgblue,       line width=1.0pt, domain=1.5707963:3.1415927] {abs(sin(deg(x)))};
  \addplot[birrenorange, line width=1.0pt, domain=3.1415927:4.712389]  {abs(cos(deg(x)))};
  \addplot[birrenred,    line width=1.0pt, domain=4.712389:5.45]       {abs(sin(deg(x)))};

  % The three roots
  \addplot[only marks, mark=*, mark size=2pt, spacecadet]
    coordinates {(1.25235,0.31309)};
  \addplot[only marks, mark=*, mark size=2pt, cgblue]
    coordinates {(2.47458,0.61864)};
  \addplot[only marks, mark=*, mark size=2pt, birrenorange]
    coordinates {(3.59530,0.89883)};

  \node[lbl, text=spacecadet,   font=\scriptsize, fill=white, inner sep=1pt]
    at (axis cs:1.25235,0.17) {$n=1$};
  \node[lbl, text=cgblue,       font=\scriptsize, fill=white, inner sep=1pt]
    at (axis cs:2.47458,0.47) {$n=2$};
  \node[lbl, text=birrenorange, font=\scriptsize, fill=white, inner sep=1pt]
    at (axis cs:3.30000,0.56) {$n=3$};

  % Curve identities, at the top of the first two arches
  \node[lbl, anchor=south west] at (axis cs:0.06,1.02) {$|\cos u|$};
  \node[lbl, anchor=south west] at (axis cs:1.63,1.02) {$|\sin u|$};

  % The straight line, named where it runs alone under the top border
  \node[lbl, anchor=east] at (axis cs:4.95,1.30) {$u/\alpha$, $\alpha=4$};

  % The excluded branches, named once, on the row above
  \node[lbl, text=birrengray, font=\scriptsize, anchor=west]
    at (axis cs:0.10,1.46)
    {\tpcifenglish{dotted: excluded branches}{gepunktet: verworfene Zweige}};
\end{axis}

\end{tikzpicture}
\end{center}


On the right the straight line $u/\alpha$ cuts only the descending arcs: for $\alpha=4$ three times; the fourth it no longer reaches. On the left the same three states in the well -- the more weakly bound, the further the tail creeps into the forbidden region.

From the solutions found for $u$ it follows that $q = \frac{2u}{L}$ and $E = \frac{\hbar^2q^2}{2m}-V_0$. We can pass to the infinitely deep potential well by taking the limit
\begin{equation}
	V_0\rightarrow \infty\quad \text{i.e.} \quad \alpha\rightarrow\infty
\end{equation}
in which we shall observe infinitely many intersection points for (\ref{eq:symmcond}) and (\ref{eq:unsymm}). These give us infinitely many bound solutions.

